% GENERATED FILE. Edit canonical lessons, then run npm run generate:books.
% canonical-source-hash: fnv1a64:096dcfe4be2626bf
% canonical-lessons: KA-C281-mayavagu, KA-C281-pyak-madu, KA-C281-kalacu, KA-C281-phelagu, KA-C281-sikkihakikollu

\chapter{To vanish, To pack, To take off, To fail, To get stuck}
\label{ch:ka-to-vanish-to-pack-to-take-off-to-fail-to-get-stuck}
\hlchaptermodality{281}

\begin{chapteropening}
\textbf{By the end of this chapter:} \emph{I can say to vanish, to pack, to take off, to fail and to get stuck.}

Complete the last lesson of chapter 281: I can say to vanish, to pack, to take off, to fail and to get stuck.
\end{chapteropening}

\section[\texorpdfstring{\kn{ಮಾಯವಾಗು} (māyavāgu)}{ಮಾಯವಾಗು (māyavāgu)}]{\kn{ಮಾಯವಾಗು} (māyavāgu) — to vanish, to disappear}

\label{lesson:KA-C281-mayavagu}



\begin{quote}
Before the new one: say the Kannada for to deliver, then the Kannada for to overtake.
\end{quote}

\subsection*{You'll want to know: \kn{ಮಾಯವಾಗು}}
\textbf{\kn{ಮಾಯವಾಗು}} — \emph{māyavāgu} — \textquotedblleft{}to vanish, to disappear\textquotedblright{}.

\begin{grammarlens}[title={What you've built}]
One more word for this chapter.
\end{grammarlens}

\subsection*{Your turn}
\begin{itemize}
\raggedright
  \item \emph{Say it:} \emph{māyavāgu}
  \item \emph{Say it:} \emph{māyavāgu}, once more
  \item \emph{Say it:} say \emph{māyavāgu}
  \item \emph{Recall:} say the Kannada for to deliver, then the Kannada for to overtake, then say \emph{māyavāgu} again
\end{itemize}

\begin{tcolorbox}[breakable,colback=teal!4,colframe=teal!35!black,title={Before you move on}]
What does \kn{ಮಾಯವಾಗು} mean? (\textquotedblleft{}To vanish, to disappear\textquotedblright{}.) Say it once more.
\end{tcolorbox}
\section[\texorpdfstring{\kn{ಪ್ಯಾಕ್} \kn{ಮಾಡು} (pyāk māḍu)}{ಪ್ಯಾಕ್ ಮಾಡು (pyāk māḍu)}]{\kn{ಪ್ಯಾಕ್} \kn{ಮಾಡು} (pyāk māḍu) — to pack}

\label{lesson:KA-C281-pyak-madu}



\begin{quote}
Before the new one: say the Kannada for to overtake, then the Kannada for to vanish.
\end{quote}

\subsection*{You'll want to know: \kn{ಪ್ಯಾಕ್} \kn{ಮಾಡು}}
\textbf{\kn{ಪ್ಯಾಕ್} \kn{ಮಾಡು}} — \emph{pyāk māḍu} — \textquotedblleft{}to pack\textquotedblright{}.

\begin{grammarlens}[title={What you've built}]
One more word for this chapter.
\end{grammarlens}

\subsection*{Your turn}
\begin{itemize}
\raggedright
  \item \emph{Say it:} \emph{pyāk māḍu}
  \item \emph{Say it:} \emph{pyāk māḍu}, once more
  \item \emph{Say it:} say \emph{pyāk māḍu}
  \item \emph{Recall:} say the Kannada for to overtake, then the Kannada for to vanish, then say \emph{pyāk māḍu} again
\end{itemize}

\begin{tcolorbox}[breakable,colback=teal!4,colframe=teal!35!black,title={Before you move on}]
What does \kn{ಪ್ಯಾಕ್} \kn{ಮಾಡು} mean? (\textquotedblleft{}To pack\textquotedblright{}.) Say it once more.
\end{tcolorbox}
\section[\texorpdfstring{\kn{ಕಳಚು} (kaḷacu)}{ಕಳಚು (kaḷacu)}]{\kn{ಕಳಚು} (kaḷacu) — to take off (clothes), to detach}

\label{lesson:KA-C281-kalacu}



\begin{quote}
Before the new one: say the Kannada for to vanish, then the Kannada for to pack.
\end{quote}

\subsection*{You'll want to know: \kn{ಕಳಚು}}
\textbf{\kn{ಕಳಚು}} — \emph{kaḷacu} — \textquotedblleft{}to take off (clothes), to detach\textquotedblright{}.

\begin{grammarlens}[title={What you've built}]
One more word for this chapter.
\end{grammarlens}

\subsection*{Your turn}
\begin{itemize}
\raggedright
  \item \emph{Say it:} \emph{kaḷacu}
  \item \emph{Say it:} \emph{kaḷacu}, once more
  \item \emph{Say it:} say \emph{kaḷacu}
  \item \emph{Recall:} say the Kannada for to vanish, then the Kannada for to pack, then say \emph{kaḷacu} again
\end{itemize}

\begin{tcolorbox}[breakable,colback=teal!4,colframe=teal!35!black,title={Before you move on}]
What does \kn{ಕಳಚು} mean? (\textquotedblleft{}To take off (clothes), to detach\textquotedblright{}.) Say it once more.
\end{tcolorbox}
\section[\texorpdfstring{\kn{ಫೇಲಾಗು} (phēlāgu)}{ಫೇಲಾಗು (phēlāgu)}]{\kn{ಫೇಲಾಗು} (phēlāgu) — to fail (an exam)}

\label{lesson:KA-C281-phelagu}



\begin{quote}
Before the new one: say the Kannada for to pack, then the Kannada for to take off.
\end{quote}

\subsection*{You'll want to know: \kn{ಫೇಲಾಗು}}
\textbf{\kn{ಫೇಲಾಗು}} — \emph{phēlāgu} — \textquotedblleft{}to fail (an exam)\textquotedblright{}.

\begin{grammarlens}[title={What you've built}]
One more word for this chapter.
\end{grammarlens}

\subsection*{Your turn}
\begin{itemize}
\raggedright
  \item \emph{Say it:} \emph{phēlāgu}
  \item \emph{Say it:} \emph{phēlāgu}, once more
  \item \emph{Say it:} say \emph{phēlāgu}
  \item \emph{Recall:} say the Kannada for to pack, then the Kannada for to take off, then say \emph{phēlāgu} again
\end{itemize}

\begin{tcolorbox}[breakable,colback=teal!4,colframe=teal!35!black,title={Before you move on}]
What does \kn{ಫೇಲಾಗು} mean? (\textquotedblleft{}To fail (an exam)\textquotedblright{}.) Say it once more.
\end{tcolorbox}
\section[\texorpdfstring{\kn{ಸಿಕ್ಕಿಹಾಕಿಕೊಳ್ಳು} (sikkihākikoḷḷu)}{ಸಿಕ್ಕಿಹಾಕಿಕೊಳ್ಳು (sikkihākikoḷḷu)}]{\kn{ಸಿಕ್ಕಿಹಾಕಿಕೊಳ್ಳು} (sikkihākikoḷḷu) — to get stuck, to get caught}

\label{lesson:KA-C281-sikkihakikollu}



\begin{quote}
Before the new one: say the Kannada for to take off, then the Kannada for to fail.
\end{quote}

\subsection*{You'll want to know: \kn{ಸಿಕ್ಕಿಹಾಕಿಕೊಳ್ಳು}}
\textbf{\kn{ಸಿಕ್ಕಿಹಾಕಿಕೊಳ್ಳು}} — \emph{sikkihākikoḷḷu} — \textquotedblleft{}to get stuck, to get caught\textquotedblright{}.

\begin{grammarlens}[title={What you've built}]
One more word for this chapter.
\end{grammarlens}

\subsection*{Your turn}
\begin{itemize}
\raggedright
  \item \emph{Say it:} \emph{sikkihākikoḷḷu}
  \item \emph{Say it:} \emph{sikkihākikoḷḷu}, once more
  \item \emph{Say it:} say \emph{sikkihākikoḷḷu}
  \item \emph{Recall:} say the Kannada for to take off, then the Kannada for to fail, then say \emph{sikkihākikoḷḷu} again
\end{itemize}

\begin{tcolorbox}[breakable,colback=teal!4,colframe=teal!35!black,title={Before you move on}]
What does \kn{ಸಿಕ್ಕಿಹಾಕಿಕೊಳ್ಳು} mean? (\textquotedblleft{}To get stuck, to get caught\textquotedblright{}.) Say it once more.
\end{tcolorbox}
